% ============================================================
% Chapter 3 — Progress Evaluation
% Drop this file's contents into main.tex (e.g. via \input{progress_evaluation}
% right after Chapter 2 / Section 2.2.1, and before \end{document}).
% Requires \usepackage{graphicx} in your preamble (main.tex already
% uses it for the diagrams in Chapter 2).
%
% Place the two screenshots in your images/figures folder and update
% the two \includegraphics paths below to match, e.g.:
%   diagrams/fig3_1_signin.png
%   diagrams/fig3_2_dashboard.png
% ============================================================

\chapter{Progress Evaluation}

\section{Current Status and Validation}

Following the design described in Chapter 2, a working frontend prototype of the
Collaborative Study Platform has been implemented and deployed locally for testing.
The prototype validates the User Login use case (UC-01, Table 2.1) and the
authentication portion of the activity, sequence, and collaboration diagrams in
Section 2.1.1: a student can register a new account or sign in with existing
credentials, and on success is taken to a personalised dashboard.

\begin{figure}[h!]
    \centering
    \includegraphics[width=0.55\textwidth]{diagrams/fig3_1_signin.png}
    \caption{Landing / authentication screen of the Collaborative Study Platform
    prototype, showing the feature highlights alongside the Sign In and Create
    Account forms.}
    \label{fig:signin-screen}
\end{figure}

\begin{figure}[h!]
    \centering
    \includegraphics[width=0.55\textwidth]{diagrams/fig3_2_dashboard.png}
    \caption{Dashboard displayed immediately after a successful registration,
    confirming the authentication flow and exposing the Notes, Join/Create Class,
    Books, and PYQs modules to the user.}
    \label{fig:dashboard-screen}
\end{figure}

As shown in Figure~\ref{fig:dashboard-screen}, registration returns a success
confirmation and routes the user directly to the Dashboard, consistent with the
post-conditions listed in Table 2.1 and with steps 6--7 of the Sequence Diagram
in Figure 2.4. The four dashboard cards correspond to the View Notes, View Books,
View Previous Year Questions, and Join/Create Class use cases identified in the
UML Use Case Diagram (Figure 2.1).

\section{Analysis of Progress}

Measured against the four project objectives set out in Section 1.2, progress to
date is as follows:

\begin{itemize}
    \item \textbf{Centralized resource access (Objective 1):} the Dashboard now
    exposes Notes, Books, PYQs, and Join/Create Class as distinct modules; the
    screens are built, and wiring them to live academic-context data is the next
    step.

    \item \textbf{Secure authentication (Objective 2):} largely complete.
    Registration and sign-in screens are functional end-to-end against the
    Go/Gin backend, and a successful sign-up correctly returns a confirmation
    and redirects to the Dashboard, as captured in Figure~\ref{fig:dashboard-screen}.

    \item \textbf{Academic-context navigation (Objective 3):} the Dashboard and
    its four resource cards are in place; the classroom $\rightarrow$
    year/stream $\rightarrow$ subject selection flow from the Swimlane Diagram
    (Figure 2.3) is designed but not yet connected to these screens.

    \item \textbf{Study-material upload (Objective 4):} not yet started on the
    frontend; the corresponding backend upload and validation endpoints
    described in Section 2.1 are planned for the next development cycle.
\end{itemize}

Overall, the authentication module and the dashboard shell are functionally
complete and match the UML use case, sequence, and collaboration diagrams
presented in Chapter 2. The remaining work is concentrated in the
resource-retrieval and upload modules, both of which are architecturally
defined but still to be implemented on the frontend.
